% ============================================================================
% MUSTERLÖSUNG - Physik Klasse 6: Weg-Zeit-Diagramme
% ============================================================================
% Identisches Layout wie AB, Lösungen in ROT
% ============================================================================

\documentclass[11pt,a4paper]{article}

\usepackage[utf8]{inputenc}
\usepackage[T1]{fontenc}
\usepackage[ngerman]{babel}
\usepackage{geometry}
\usepackage{helvet}
\usepackage{tikz}
\usepackage{amsmath}
\usepackage{amssymb}
\usepackage{enumitem}
\usepackage{tabularx}
\usepackage{array}
\usepackage{fancyhdr}
\usepackage{lastpage}
\usepackage{xcolor}

\renewcommand{\familydefault}{\sfdefault}

\geometry{left=2cm, right=2cm, top=2.5cm, bottom=2cm}

% Farben
\definecolor{physik}{HTML}{2c5aa0}
\definecolor{lightblue}{HTML}{e3f2fd}
\definecolor{lightgray}{HTML}{f5f5f5}
\definecolor{loesung}{HTML}{c62828}

% Kopfzeile
\pagestyle{fancy}
\fancyhf{}
\fancyhead[L]{\small\textcolor{loesung}{\textbf{MUSTERLÖSUNG}} | Physik Klasse 6 | Weg-Zeit-Diagramme}
\fancyhead[R]{\small Name: \underline{\hspace{4cm}}}
\fancyfoot[C]{\small Seite \thepage{} von \pageref{LastPage}}
\renewcommand{\headrulewidth}{0.4pt}

% Lösung-Befehl
\newcommand{\lsg}[1]{\textcolor{loesung}{\textbf{#1}}}

% Kasten für Übertragung
\newcommand{\uebertragkasten}[2]{%
    \vspace{0.5em}
    \noindent\fcolorbox{physik}{lightblue}{%
        \begin{minipage}{\dimexpr\textwidth-2\fboxsep-2\fboxrule}
            \textbf{Kasten #1: #2}\\[0.8em]
            \rule{0pt}{2.5em}
        \end{minipage}%
    }%
    \vspace{0.5em}
}

\newcommand{\uebertragkastenGross}[2]{%
    \vspace{0.5em}
    \noindent\fcolorbox{physik}{lightblue}{%
        \begin{minipage}{\dimexpr\textwidth-2\fboxsep-2\fboxrule}
            \textbf{Kasten #1: #2}\\[0.8em]
            \rule{0pt}{5em}
        \end{minipage}%
    }%
    \vspace{0.5em}
}

% Lücke mit Lösung
\newcommand{\lueckelsg}[1]{\underline{\lsg{\,#1\,}}}

% Niveaustufen
\newcommand{\stern}[1]{\textbf{(#1)}}

\begin{document}

% ============================================================================
% TITEL
% ============================================================================
\begin{center}
    {\Large\bfseries\textcolor{physik}{Weg-Zeit-Diagramme verstehen und zeichnen}}\\[0.3em]
    {\normalsize Physik Klasse 6 | \textcolor{loesung}{\textbf{MUSTERLÖSUNG}}}
\end{center}

\vspace{0.5em}
\hrule
\vspace{1em}

% ============================================================================
% KASTEN 1: Achsenbeschriftung
% ============================================================================
\uebertragkasten{1}{Achsenbeschriftung im s(t)-Diagramm}

\noindent Ergänze die Beschriftung:

\begin{itemize}[leftmargin=2em]
    \item Waagerechte Achse (x-Achse): \lueckelsg{Zeit $t$} in \lueckelsg{Sekunden (s)}
    \item Senkrechte Achse (y-Achse): \lueckelsg{Weg $s$} in \lueckelsg{Metern (m)}
\end{itemize}

\vspace{1em}

% ============================================================================
% KASTEN 2: Merksatz
% ============================================================================
\uebertragkastenGross{2}{Interpretation des Diagramms}

\noindent Ergänze den Merksatz:

\begin{itemize}[leftmargin=2em]
    \item Eine \textbf{steile Gerade} bedeutet: \lueckelsg{schnelle} Bewegung
    \item Eine \textbf{flache Gerade} bedeutet: \lueckelsg{langsame} Bewegung
    \item Eine \textbf{waagerechte Linie} bedeutet: \lueckelsg{Stillstand}
\end{itemize}

\vspace{1.5em}

% ============================================================================
% AUFGABE 1: Werte ablesen (*)
% ============================================================================
\noindent\textbf{Aufgabe 1 \stern{*}: Werte aus dem Diagramm ablesen} \hfill \textit{/6 P}

\vspace{0.5em}
\noindent Das Diagramm zeigt die Bewegung eines Radfahrers:

\begin{center}
\begin{tikzpicture}[scale=0.85]
    % Koordinatensystem
    \draw[->] (0,0) -- (8.5,0) node[right] {$t$ in s};
    \draw[->] (0,0) -- (0,5.5) node[above] {$s$ in m};

    % Skalierung x
    \foreach \x/\label in {0/0, 1/2, 2/4, 3/6, 4/8, 5/10} {
        \draw (\x,-0.1) -- (\x,0.1);
        \node[below] at (\x,-0.2) {\small \label};
    }

    % Skalierung y
    \foreach \y/\label in {0/0, 1/20, 2/40, 3/60, 4/80, 5/100} {
        \draw (-0.1,\y) -- (0.1,\y);
        \node[left] at (-0.2,\y) {\small \label};
    }

    % Gerade (Radfahrer: v = 10 m/s)
    \draw[thick, physik] (0,0) -- (5,5);

    % Lösungs-Markierungen
    \draw[dashed, loesung] (2,0) -- (2,2) -- (0,2);
    \fill[loesung] (2,2) circle (3pt);
    \node[loesung, right] at (2.1,2) {\small (4|40)};

    % Beschriftung
    \node[physik, right] at (4.5,4.5) {\small Radfahrer};
\end{tikzpicture}
\end{center}

\noindent Lies die Werte ab und ergänze:

\begin{enumerate}[label=\alph*), leftmargin=2em]
    \item Nach $t = 4$ s hat der Radfahrer $s =$ \lueckelsg{40 m} zurückgelegt. \hfill \lsg{2 P}
    \item Nach $t = 10$ s hat der Radfahrer $s =$ \lueckelsg{100 m} zurückgelegt. \hfill \lsg{2 P}
    \item Für eine Strecke von $s = 60$ m benötigt der Radfahrer $t =$ \lueckelsg{6 s}. \hfill \lsg{2 P}
\end{enumerate}

\vspace{1.5em}

% ============================================================================
% AUFGABE 2: Diagramm zeichnen (**)
% ============================================================================
\noindent\textbf{Aufgabe 2 \stern{**}: Ein Diagramm zeichnen} \hfill \textit{/8 P}

\vspace{0.5em}
\noindent Ein Jogger läuft mit gleichförmiger Geschwindigkeit. Die Messwerte:

\begin{center}
\begin{tabular}{|c|c|c|c|c|c|}
\hline
Zeit $t$ in s & 0 & 5 & 10 & 15 & 20 \\
\hline
Weg $s$ in m & 0 & 25 & 50 & 75 & 100 \\
\hline
\end{tabular}
\end{center}

\noindent\textbf{a)} Zeichne die Werte als Punkte in das Koordinatensystem ein.

\noindent\textbf{b)} Verbinde die Punkte mit dem Lineal zu einer Geraden.

\begin{center}
\begin{tikzpicture}[scale=0.7]
    % Koordinatensystem
    \draw[->] (0,0) -- (11,0) node[right] {$t$ in s};
    \draw[->] (0,0) -- (0,6) node[above] {$s$ in m};

    % Skalierung x
    \foreach \x/\label in {0/0, 2.5/5, 5/10, 7.5/15, 10/20} {
        \draw (\x,-0.1) -- (\x,0.1);
        \node[below] at (\x,-0.3) {\small \label};
    }

    % Skalierung y
    \foreach \y/\label in {0/0, 1.25/25, 2.5/50, 3.75/75, 5/100} {
        \draw (-0.1,\y) -- (0.1,\y);
        \node[left] at (-0.3,\y) {\small \label};
    }

    % Hilfslinien
    \draw[lightgray, very thin] (0,1.25) -- (10,1.25);
    \draw[lightgray, very thin] (0,2.5) -- (10,2.5);
    \draw[lightgray, very thin] (0,3.75) -- (10,3.75);
    \draw[lightgray, very thin] (0,5) -- (10,5);
    \draw[lightgray, very thin] (2.5,0) -- (2.5,5);
    \draw[lightgray, very thin] (5,0) -- (5,5);
    \draw[lightgray, very thin] (7.5,0) -- (7.5,5);
    \draw[lightgray, very thin] (10,0) -- (10,5);

    % LÖSUNG: Punkte und Gerade
    \draw[thick, loesung] (0,0) -- (10,5);
    \fill[loesung] (0,0) circle (4pt);
    \fill[loesung] (2.5,1.25) circle (4pt);
    \fill[loesung] (5,2.5) circle (4pt);
    \fill[loesung] (7.5,3.75) circle (4pt);
    \fill[loesung] (10,5) circle (4pt);

    \node[loesung, right] at (9,4.2) {\small Jogger};
\end{tikzpicture}
\end{center}

\lsg{Punkte korrekt: 5 P | Gerade korrekt: 1 P}

\noindent\textbf{c)} Berechne die Geschwindigkeit des Joggers:

\vspace{0.3em}
\hspace{2em} $v = \dfrac{s}{t} = \dfrac{\lueckelsg{100 m}}{\lueckelsg{20 s}} =$ \lueckelsg{5 m/s} \hfill \lsg{2 P}

\vspace{1.5em}

% ============================================================================
% AUFGABE 3: Bewegungen vergleichen (**)
% ============================================================================
\noindent\textbf{Aufgabe 3 \stern{**}: Bewegungen vergleichen} \hfill \textit{/6 P}

\vspace{0.5em}
\noindent Das Diagramm zeigt zwei Bewegungen:

\begin{center}
\begin{tikzpicture}[scale=0.7]
    % Koordinatensystem
    \draw[->] (0,0) -- (8.5,0) node[right] {$t$ in s};
    \draw[->] (0,0) -- (0,5.5) node[above] {$s$ in m};

    % Skalierung x
    \foreach \x/\label in {0/0, 2/4, 4/8, 6/12, 8/16} {
        \draw (\x,-0.1) -- (\x,0.1);
        \node[below] at (\x,-0.2) {\small \label};
    }

    % Skalierung y
    \foreach \y/\label in {0/0, 1/20, 2/40, 3/60, 4/80} {
        \draw (-0.1,\y) -- (0.1,\y);
        \node[left] at (-0.2,\y) {\small \label};
    }

    % Person A (schneller)
    \draw[thick, red] (0,0) -- (4,4);
    \node[red, right] at (3.5,3.8) {\small A};

    % Person B (langsamer)
    \draw[thick, physik] (0,0) -- (8,4);
    \node[physik, right] at (7,3.5) {\small B};

    % Lösungs-Markierung
    \draw[dashed, loesung] (4,0) -- (4,4);
    \node[loesung] at (4,-0.6) {\small 8 s};
    \node[loesung, left] at (4,4) {\small 80 m};
\end{tikzpicture}
\end{center}

\begin{enumerate}[label=\alph*), leftmargin=2em]
    \item Welche Person ist schneller? \lueckelsg{Person A} \hfill \lsg{1 P}
    \item Begründe deine Antwort:\\[0.5em]
          \lsg{Die Gerade von Person A ist steiler als die von Person B.}\\[0.5em]
          \lsg{Eine steilere Gerade bedeutet eine höhere Geschwindigkeit.} \hfill \lsg{3 P}
    \item Wie weit ist Person A nach 8 Sekunden gekommen? $s =$ \lueckelsg{80 m} \hfill \lsg{2 P}
\end{enumerate}

\vspace{1.5em}

% ============================================================================
% AUFGABE 4: Bewegungsgeschichte (***)
% ============================================================================
\noindent\textbf{Aufgabe 4 \stern{***}: Eine Bewegungsgeschichte erfinden} \hfill \textit{/5 P}

\vspace{0.5em}
\noindent Erfinde zu dem folgenden Diagramm eine passende Geschichte:

\begin{center}
\begin{tikzpicture}[scale=0.65]
    % Koordinatensystem
    \draw[->] (0,0) -- (11,0) node[right] {$t$ in min};
    \draw[->] (0,0) -- (0,5.5) node[above] {$s$ in m};

    % Skalierung x
    \foreach \x/\label in {0/0, 2/2, 4/4, 6/6, 8/8, 10/10} {
        \draw (\x,-0.1) -- (\x,0.1);
        \node[below] at (\x,-0.2) {\small \label};
    }

    % Skalierung y
    \foreach \y/\label in {0/0, 1/100, 2/200, 3/300, 4/400, 5/500} {
        \draw (-0.1,\y) -- (0.1,\y);
        \node[left] at (-0.2,\y) {\small \label};
    }

    % Bewegung
    \draw[thick, physik] (0,0) -- (3,3);
    \draw[thick, physik] (3,3) -- (6,3);
    \draw[thick, physik] (6,3) -- (10,5);

    % Beschriftungen für Lösung
    \node[loesung, above] at (1.5,1.5) {\small schnell};
    \node[loesung, above] at (4.5,3) {\small Stillstand};
    \node[loesung, above] at (8,4) {\small langsam};
\end{tikzpicture}
\end{center}

\noindent\textbf{Meine Bewegungsgeschichte:}

\vspace{0.5em}
\noindent\lsg{Beispiel: Lisa fährt mit dem Fahrrad zur Schule. In den ersten 3 Minuten}\\[0.5em]
\noindent\lsg{fährt sie schnell und legt 300 m zurück. Dann muss sie an einer}\\[0.5em]
\noindent\lsg{roten Ampel warten (3 Minuten Stillstand). Danach fährt sie}\\[0.5em]
\noindent\lsg{langsamer weiter und erreicht nach weiteren 4 Minuten die Schule (500 m).}

\vspace{0.5em}
\lsg{Bewertung: Geschichte passt zum Diagramm: 3 P | Drei Phasen erkannt: 2 P}

\vspace{1.5em}

% ============================================================================
% PUNKTE
% ============================================================================
\hrule
\vspace{0.5em}

\begin{center}
\begin{tabular}{|c|c|c|c|}
\hline
\textbf{Niveau} & \textbf{Aufgaben} & \textbf{max. Punkte} & \textbf{erreicht} \\
\hline
\stern{*} Basis & Kästen 1-2, Aufgabe 1 & 10 P & \hspace{1.5cm}/10 \\
\hline
\stern{**} Standard & + Aufgabe 2, 3 & 24 P & \hspace{1.5cm}/24 \\
\hline
\stern{***} Erweitert & + Aufgabe 4 & 29 P & \hspace{1.5cm}/29 \\
\hline
\end{tabular}
\end{center}

\end{document}
